\let\oldprintchaptertitle=\printchaptertitle
\renewcommand{\printchaptertitle}[1]{%
	\vspace*{-75pt}
	\oldprintchaptertitle{#1}
}%
\myChapterStar{Abstract}{}{section}
\let\printchaptertitle=\oldprintchaptertitle

In a constantly evolving technological landscape, Service-Oriented Architectures (SOA) and microservices have become essential paradigms for designing modular, flexible, and distributed software systems. 
These architectures enable the decomposition of application functionalities into autonomous services, thereby promoting reusability, scalability, and adaptability. However, the way these services are composed remains a critical challenge for ensuring performance, robustness, and responsiveness—particularly in dynamic environments such as cloud computing or data-driven systems.
Traditional composition approaches —\textit{whether static or dynamic}— %face limitations when it comes to scalability, fault tolerance, and lazy evaluation. 
show limitations when confronted with modern requirements for scalability, resilience, and lazy evaluation.
% In this context, incremental service composition emerges as a promising alternative. This data-driven approach allows the execution of a process to progress step by step, based on the availability of local data, without relying on a predefined global control flow. It is notably supported by the Guarded Attribute Grammars (GAG) model, which represents service coordination as production rules driven by explicit dependencies among attributes.
% In this thesis, we propose to evaluate the performance of an incremental composition engine through a benchmark based on the Fourier Transform (FT). The Fast Fourier Transform (FFT) algorithm, in particular, is an ideal case study due to its recursive structure, natural parallelism, and the fine granularity of its operations, which can be modeled as a sequence of interdependent services.
This is where incremental service composition comes into play: a data-driven emerging approach in which the execution of a process unfolds progressively, guided by the availability of computed and expected data between services. It relies on the Guarded Attribute Grammars (GAG) model, which represents service composition logic through rules driven by explicit data dependencies. A key advantage of incremental composition is its native compatibility with containerization, a widely adopted practice in cloud architectures by major companies such as Google, Amazon, and IBM. This sets it apart from traditional approaches based on ad hoc models that do not support container-based deployments.
%------->
% In order to rigorously evaluate this approach, we proposed a benchmark dedicated to incremental Fourier transforms, deployed on a Kubernetes cluster of a physical machine.
% Experimental results demonstrate that incremental composition offers improved adaptability to execution environments, better exploitation of parallelism, and enhanced robustness against partial failures, making it a promising approach for modern distributed systems.
%---->
% This study paves the way for new avenues in the optimization of lazy service-based systems and underscores the importance of a close coupling between formal modeling (GAG) and modern infrastructures (Kubernetes).
% This work highlights the benefits of combining formal models such as GAG with cloud-native technologies like Docker, Kubernetes, and Vagrant to build efficient, intelligent, and adaptive service composition engines.
To rigorously evaluate this approach, we designed a benchmark based on incremental Fourier Transform, deployed on a Kubernetes cluster running on a physical machine. Experimental results demonstrate that incremental composition provides better adaptability to the execution environment, naturally supports parallelism to significantly reduce runtime, and enables more robust implementations capable of handling larger data volumes while offering increased fault tolerance. Our findings generalize to any algorithm that follows the Divide and Conquer paradigm.

\vspace{1cm}
\noindent\textbf{Keywords:} Service Oriented Architecture (SOA), Guarded Attribute Grammars (GAG), Lazy services, incremental composition, benchmark, Fourier Transform, Kubernetes, cloud computing.





\myCleanStarChapterEnd
